\documentclass[12pt]{article}
\usepackage[T1]{fontenc}
\usepackage[utf8]{inputenc}
\usepackage{lmodern}
\usepackage{textcomp}
\usepackage{enumitem}
\usepackage{geometry}
\geometry{margin=1in}
\begin{document}
\begin{center}
Answer Key - Political Science - Graduate Level Assessment 5 \
\small (Answers are ordered exactly as in the question paper.)
\end{center}

\section*{Multiple Choice}
\begin{enumerate}[label=\arabic*.]
\item \textbf{Answer:} B
\\ \textit{Options:} A) US expansion was based on territory rather than markets post-1865; B) US expansion was based on markets rather than territory post-1865; C) US expansion was limited to Latin America post-1865; D) US expansion ended after 1865
\\ \textit{Note:} The correct answer highlights that after 1865, following the Civil War, U.S. expansion increasingly focused on economic growth and market development, driven by industrialization and capitalism, rather than the earlier emphasis on territorial acquisition and Manifest Destiny.
\item \textbf{Answer:} B
\\ \textit{Options:} A) Political policies designed to alter demographic characteristics of a state.; B) All of these options.; C) Legislation that criminalises certain cultural behaviours or practice.; D) The use of military force to conduct ethnic-cleansing through displacement and killing.
\\ \textit{Note:} The correct answer "All of these options" is valid because societal security encompasses a range of threats, including economic instability, social unrest, and environmental crises, all of which can undermine the stability and safety of a society.
\item \textbf{Answer:} C
\\ \textit{Options:} A) To reward recipients for the cessation of negative behaviour, such as human rights abuses.; B) To incentivise recipients to act in ways beneficial to the donor.; C) Compassion for the human suffering of others.; D) To influence recipients or potential recipients through granting or denying aid.
\\ \textit{Note:} The correct answer is "Compassion for the human suffering of others" because economic aid as a policy for security primarily focuses on strategic interests and national security objectives rather than humanitarian concerns.
\item \textbf{Answer:} D
\\ \textit{Options:} A) the president of the United States.; B) the secretary of defense.; C) the chairman of the Joint Chiefs of Staff.; D) Congress.
\\ \textit{Note:} The power to declare war rests with Congress as stated in Article I, Section 8 of the U.S. Constitution, which grants Congress the authority to make decisions regarding war and military engagement, thereby ensuring a system of checks and balances between the legislative and executive branches.
\item \textbf{Answer:} A
\\ \textit{Options:} A) Withdrawal of economic trade rights with the domestic market.; B) Export controls protecting technological advantage and further foreign policy objectives.; C) Control of munitions and arms sales.; D) Import restrictions to protect a domestic market from foreign goods.
\\ \textit{Note:} The withdrawal of economic trade rights with the domestic market is the odd one out because it directly restricts a country's own economic interactions, whereas the other options typically involve leveraging economic tools to influence or sanction external entities or adversaries.
\end{enumerate}

\section*{True / False}
\begin{enumerate}[label=\arabic*.]
\setcounter{enumi}{5}
\item \textbf{Answer:} True
\\ \textit{Options:} A) True; B) False
\\ \textit{Note:} The Clinton Administration supported the transition of the Russian economy from a state-controlled system to a market-oriented one by providing economic assistance, promoting privatization, and advocating for reforms that encouraged foreign investment and entrepreneurship.
\item \textbf{Answer:} False
\\ \textit{Options:} A) True; B) False
\\ \textit{Note:} The statement is false because American unilateralism is often justified by a belief in national sovereignty and effectiveness in achieving U.S. interests, rather than relying on the success of multilateral approaches, which are frequently viewed as slow and inefficient.
\item \textbf{Answer:} True
\\ \textit{Options:} A) True; B) False
\\ \textit{Note:} The Central Intelligence Agency (CIA) is mandated by law to collect, analyze, and disseminate intelligence on foreign nations to inform and guide U.S. government policymakers in matters of national security and foreign relations.
\item \textbf{Answer:} True
\\ \textit{Options:} A) True; B) False
\\ \textit{Note:} The statement is true because the United Nations Security Council is structured with five permanent members-China, France, Russia, the United Kingdom, and the United States-who possess veto power over resolutions, while the other ten non-permanent members are elected for two-year terms and do not have veto privileges.
\item \textbf{Answer:} True
\\ \textit{Options:} A) True; B) False
\\ \textit{Note:} American exceptionalism was enhanced during the Cold War as the United States positioned itself as the foremost defender of democracy and capitalism against communism, promoting the belief in its unique role and values on the global stage.
\end{enumerate}

\section*{Long Form Response}
\begin{enumerate}[label=\arabic*.]
\setcounter{enumi}{10}
\item \textbf{Answer:} Bill Clinton's assertion that "Globalization is not something we can hold off or turn off" underscores the inevitability of interconnectedness in today's political and economic landscape. This perspective implies that nations must adapt to globalization rather than resist it, as attempts to isolate economies can lead to greater instability and inequity. Contemporary challenges, including rising protectionism, economic inequality, and the impact of technological advancements, highlight the complexities of globalization that require cooperative international responses. Clinton's statement serves as a reminder that political leaders must navigate these challenges by fostering policies that embrace global interdependence while addressing domestic concerns. Ultimately, it suggests that proactive engagement with globalization is essential for sustainable development and political stability in an increasingly interconnected world.
\\ \textit{Note:} The answer correctly emphasizes that Clinton's assertion reflects the necessity for countries to embrace globalization in order to effectively address contemporary issues such as protectionism and economic inequality, as resisting these global trends can exacerbate instability.
\item \textbf{Answer:} Global and regional international trade agreements operate through various mechanisms, including tariff reductions, trade facilitation measures, and regulatory harmonization, which collectively shape the economic landscape between nations. These agreements often foster deeper political relations by encouraging cooperation and dialogue, as countries become economically interdependent. Additionally, they can lead to power dynamics where stronger nations influence the terms to their advantage, potentially marginalizing weaker economies. The implications of these agreements extend beyond economics, as they can shift alliances and alter geopolitical strategies, highlighting the intricate interplay between trade and international relations. Overall, the multifaceted nature of these agreements underscores their significance in shaping both political and economic ties among nations.
\\ \textit{Note:} The answer correctly identifies key mechanisms of international trade agreements-such as tariff reductions and regulatory harmonization-that not only enhance economic interdependence among nations but also facilitate political cooperation, while acknowledging the potential for power imbalances in negotiations.
\end{enumerate}

\vfill
\noindent\textit{End of Answer Key}
\end{document}